\documentclass[11pt]{article}
\pagestyle{plain}

% \pgfplotsset{compat=1.18}
\usepackage[margin=.9in,top=.8in]{geometry}
\usepackage{amsmath,amsfonts,amssymb,amsthm}
\usepackage[ruled,vlined,linesnumbered,commentsnumbered]{algorithm2e}
\SetKwInput{KwIn}{Input}
\SetKwInput{KwOut}{Output}
\SetKwInput{KwGlobal}{Global variable}
\SetKwInput{KwOracle}{Oracle}
\SetKwComment{Comment}{// }{}
\usepackage{graphicx}
\usepackage{enumerate}
\usepackage{bbm}
\usepackage{verbatim}
\usepackage{hyperref,color}
\usepackage[capitalize,nameinlink]{cleveref}
\usepackage[dvipsnames]{xcolor}
\hypersetup{
	colorlinks=true,
	pdfpagemode=UseNone,
	citecolor=OliveGreen,
	linkcolor=NavyBlue,
	urlcolor=Magenta,
	pdfstartview=FitW
}
\usepackage{appendix}
\usepackage{makecell}
\usepackage{tablefootnote}
\crefname{appsec}{Appendix}{Appendices}
\usepackage{tikz}
\usepackage{pgfplots}
\usepackage{xifthen}
\usepackage{subcaption}
\usepackage{hhline}
\usepackage{listings}

\theoremstyle{plain}
\newtheorem{theorem}{Theorem}[section]
\newtheorem{proposition}[theorem]{Proposition}
\newtheorem{lemma}[theorem]{Lemma}
\newtheorem{corollary}[theorem]{Corollary}
\newtheorem{conjecture}[theorem]{Conjecture}
\newtheorem{problem}[theorem]{Problem}
\newtheorem{claim}[theorem]{Claim}
\newtheorem{fact}[theorem]{Fact}

\theoremstyle{definition}
\newtheorem{definition}[theorem]{Definition}
\newtheorem{example}[theorem]{Example}
\newtheorem{observation}[theorem]{Observation}
\newtheorem{assumption}[theorem]{Assumption}
\newtheorem*{assumption*}{Assumption}
\newtheorem{condition}[theorem]{Condition}

\newtheorem*{prob}{Proposed Problem}

\theoremstyle{remark}
\newtheorem{remark}[theorem]{Remark}

\crefname{lemma}{Lemma}{Lemmas}
\crefname{theorem}{Theorem}{Theorems}
\crefname{definition}{Definition}{Definitions}
\crefname{fact}{Fact}{Facts}
\crefname{claim}{Claim}{Claims}
\crefname{proposition}{Proposition}{Propositions}


\newcommand{\dif}{\,\mathrm{d}}
%\newcommand{\E}{\mathbb{E}}
%\newcommand{\Var}{\mathrm{Var}}
\newcommand{\Cov}{\mathrm{Cov}}
\newcommand{\MI}{\mathrm{I}}
%\newcommand{\Ent}{\mathrm{Ent}}
\newcommand{\ind}{\mathbbm{1}}
\newcommand{\allone}{\mathbf{1}}
\newcommand{\tr}{\mathrm{tr}}
\newcommand{\wt}{\mathrm{weight}}
\DeclareMathOperator*{\argmax}{arg\,max}
\newcommand{\norm}[1]{\left\lVert #1 \right\rVert}
\newcommand{\TV}[2]{\left\|#1 - #2\right\|_{\mathrm{TV}}}
\newcommand{\ceil}[1]{\left\lceil #1 \right\rceil}
\newcommand{\floor}[1]{\left\lfloor #1 \right\rfloor}
\newcommand{\ip}[2]{\left\langle #1 , #2 \right\rangle}
\newcommand{\grad}{\nabla}
\newcommand{\hessian}{\nabla^2}
\newcommand{\trans}{\intercal}
\newcommand{\diag}{\mathrm{diag}}
\newcommand{\poly}{\mathrm{poly}}
\newcommand{\dist}{\mathrm{dist}}
\newcommand{\sgn}{\mathrm{sgn}}
\newcommand{\fpras}{\mathsf{FPRAS}}
\newcommand{\fpaus}{\mathsf{FPAUS}}
\newcommand{\fptas}{\mathsf{FPTAS}}

\newcommand{\eps}{\varepsilon}

\newcommand{\N}{\mathbb{N}}
\newcommand{\Z}{\mathbb{Z}}
\newcommand{\Q}{\mathbb{Q}}
\newcommand{\R}{\mathbb{R}}
\newcommand{\C}{\mathbb{C}}

\renewcommand{\AA}{\mathcal{A}}
\newcommand{\BB}{\mathcal{B}}
\newcommand{\CC}{\mathcal{C}}
\newcommand{\DD}{\mathcal{D}}
\newcommand{\EE}{\mathcal{E}}
\newcommand{\FF}{\mathcal{F}}
\newcommand{\GG}{\mathcal{G}}
\newcommand{\HH}{\mathcal{H}}
\newcommand{\II}{\mathcal{I}}
\newcommand{\JJ}{\mathcal{J}}
\newcommand{\KK}{\mathcal{K}}
\newcommand{\LL}{\mathcal{L}}
\newcommand{\MM}{\mathcal{M}}
\newcommand{\NN}{\mathcal{N}}
\newcommand{\OO}{\mathcal{O}}
\newcommand{\PP}{\mathcal{P}}
\newcommand{\QQ}{\mathcal{Q}}
\newcommand{\RR}{\mathcal{R}}
\renewcommand{\SS}{\mathcal{S}}
\newcommand{\TT}{\mathcal{T}}
\newcommand{\UU}{\mathcal{U}}
\newcommand{\VV}{\mathcal{V}}
\newcommand{\WW}{\mathcal{W}}
\newcommand{\XX}{\mathcal{X}}
\newcommand{\YY}{\mathcal{Y}}
\newcommand{\ZZ}{\mathcal{Z}}

\newcommand{\KL}[2]{D_{\mathrm{KL}} \left( #1 \,\Vert\, #2 \right)}
\newcommand{\KLbig}[2]{D_{\mathrm{KL}} \left( #1 \,\big\Vert\, #2 \right)}
\newcommand{\KLBig}[2]{D_{\mathrm{KL}} \left( #1 \,\Big\Vert\, #2 \right)}

\newcommand{\SKL}[2]{D_{\mathrm{SKL}} \left( #1 \,\Vert\, #2 \right)}
\newcommand{\SKLbig}[2]{D_{\mathrm{SKL}} \left( #1 \,\big\Vert\, #2 \right)}
\newcommand{\SKLBig}[2]{D_{\mathrm{SKL}} \left( #1 \,\Big\Vert\, #2 \right)}

\newcommand{\qandq}{\quad\text{and}\quad}
\newcommand{\supp}{\mathrm{supp}}

\newcommand{\DKL}[2]{\-D_{\-{KL}}\left(#1 \parallel #2\right)}
\newcommand{\DTV}[2]{\-D_{\mathrm{TV}}\left({#1},{#2}\right)}
% \newcommand{\dist}{\mathrm{dist}}
\newcommand{\e}{\mathrm{e}}
\renewcommand{\epsilon}{\varepsilon}
\renewcommand{\emptyset}{\varnothing}
% \newcommand{\norm}[1]{\left\Vert#1\right\Vert}
\newcommand{\set}[1]{\left\{#1\right\}}
\newcommand{\tuple}[1]{\left(#1\right)} 
% \newcommand{\eps}{\varepsilon}
\newcommand{\inner}[2]{\left\langle #1,#2\right\rangle}
\newcommand{\tp}{\tuple}
\newcommand{\ol}{\overline}
\newcommand{\abs}[1]{\left\vert#1\right\vert}
\newcommand{\ctp}[1]{\left\lceil#1\right\rceil}
\newcommand{\ftp}[1]{\left\lfloor#1\right\rfloor}
\newcommand{\Ind}[1]{\-{Ind}_{#1}}
\newcommand{\todo}[1]{{\color{red} TODO: #1}}

\newcommand{\Cor}{\Psi^{\mathrm{sym}}}
\newcommand{\cor}{\Psi^{\-{cor}}}
\newcommand{\Inf}{\Psi}
% \newcommand{\diag}{\-{diag}}

\def\*#1{\boldsymbol{#1}} % Use \*A for \mathbf{A}
\def\+#1{\mathcal{#1}} % Use \+A for \mathcal{A}
\def\-#1{\mathrm{#1}} % Use \-A for \mathrm{A}
\def\=#1{\mathbb{#1}} % Use \^A for \mathbb{A}
\def\!#1{\mathfrak{#1}} % Use \!A for \mathfrak{A}

\newcommand*\diff{\mathop{}\!\mathrm{d}}

% probability
% \DeclareMathOperator*{\oPr}{\mathbf{Pr}}
\def\oPr{\mathop{\mathrm{Pr}}}
\renewcommand{\Pr}[2][]{ \ifthenelse{\isempty{#1}}
  {\oPr\left[#2\right]}
  {\oPr_{#1}\left[#2\right]} } % Use \Pr[a]{b} for \mathbf{Pr}_a[b], \Pr{b} for  \mathbf{Pr}[b]

% expectation: relies on the package "dsfont"
% \DeclareMathOperator*{\oE}{\mathbb{E}}
\def\oE{\mathop{\mathbb{E}}}
\newcommand{\E}[2][]{ \ifthenelse{\isempty{#1}}
  {\oE\left[#2\right]}
  {\oE_{#1}\left[#2\right]} }

% variance
%\DeclareMathOperator*{\oVar}{\mathbf{Var}}
\def\oVar{\mathrm{Var}}
\newcommand{\Var}[2][]{ \ifthenelse{\isempty{#1}}
  {\oVar\left[#2\right]}
  {\oVar_{#1}\left[#2\right]} }

% entropy
% \DeclareMathOperator*{\oEnt}{\mathbf{Ent}}
\def\oEnt{\mathrm{Ent}}
\newcommand{\Ent}[2][]{ \ifthenelse{\isempty{#1}}
  {\oEnt\left[#2\right]}
  {\oEnt_{#1}\left[#2\right]} }

\newcommand{\PhiEnt}[2][]{ \ifthenelse{\isempty{#1}}
  {\oEnt^\phi\left[#2\right]}
  {\oEnt^\phi_{#1}\left[#2\right]} }


\newenvironment{Exampleparagraph}[1][]
{
    \par 
    \noindent 
    \textbf{\color{blue} #1}
    \par 
    \vspace{0.5em} 
    \itshape 
}
{
    \par 
    \vspace{1em} 
}

%\input{local-macro}

\newcommand{\RED}[1]{\textcolor{red}{#1}}

\newcommand{\bern}[1]{\mathrm{Bern}\left(#1\right)}
\newcommand{\pll}{\parallel}
\usepackage{titling}
\setlength{\droptitle}{-7em}


\title{Influence and Isoperimetric Inequalities on Matroids}
\date{}
\author{Zejia Chen}

\pgfplotsset{compat=1.18} 
\begin{document}

\maketitle
\vspace{-4em}
\input{intro.tex}
\paragraph{Previous Work and Proposed Problem.}
% Talagrand's $L^1$--$L^2$ inequality was first proved for biased product distribution 
% It was later extended to under suitable hypercontractivity assumptions \cite{KoehlerLifshitzMinzerMossel23}.
% Eldan and Gross proved their inequality for the uniform measure using pathwise stochastic analysis \cite{EldanGross22}.
% Ivanisvili and Zhang subsequently gave a simpler semigroup proof that also applies to biased product measures \cite{IvanisviliZhang26}.

Most previous work establishes influence and isoperimetric inequalities for uniform or biased product distributions \cite{Talagrand93,KKL88,CorderoErausquinLedoux12}.
KKL and other influence-type inequalities have also been extended to certain Schreier graphs \cite{ODonnellWimmer13} and Markov random fields on bounded-degree graphs under suitable mixing assumptions \cite{KoehlerLifshitzMinzerMossel23}. 
In my recent work with Zongchen Chen and Xinyuan Zhang \cite{chen2026discrete}, we established these inequalities for distributions with full support that satisfy marginal boundedness and a Dobrushin-type condition \cite{Dobrushin68}, which controls dependence among coordinates.

A limitation of our hypercube argument is its reliance on full support.
It therefore does not directly apply to distributions with constrained support.
An important exception is the uniform distribution over matroid independent sets.
We say $\+M=(S,\+I)$ is a matroid if $\+I\subseteq 2^S$ is downward-closed and satisfies the exchange property: $\forall A,B\in \+I, |A|>|B|\implies \exists u\in A\setminus B, B\cup\{u\}\in\+I$.
The members of $\mathcal{I}$ are its independent sets.
Matroids capture a fundamental combinatorial structure in theoretical computer science.
For example, graphic matroids describe acyclic subgraphs.
Studying influence and isoperimetric inequalities on matroid independent sets provides insights into Boolean function analysis beyond the Boolean hypercube \cite{Hop26}.
In the graphic case, isoperimetric inequalities on graphic matroids would help quantify the stability of graph properties over acyclic subgraphs.
This motivates us to explore the influence and isoperimetric inequalities on matroid independent sets.
Formally, we would like to answer the following problem.
\begin{prob}\emph{
    Do the KKL inequality and Talagrand's inequality hold for the uniform distribution over matroid independent sets?
}\end{prob}
% The hard-core model, a distribution over vertex-weighted independent sets, provides another natural setting for these questions.
% It can model selection in a social network where no two committee members are friends.
More broadly, the proposed problem is a first step toward analogous inequalities for high-dimensional expanders \cite{gotlib2023no}, determinantal point processes \cite{taskar2013determinantal}, and spin systems with hard constraints \cite{chen2023spectral}, such as the hard-core model (vertex-weighted independent sets) and graph colorings.
Further targets include Talagrand's $L^1$--$L^2$ inequality \cite{Talagrand94}, which strengthens KKL, and the Eldan--Gross inequality \cite{EldanGross22}, which incorporates both KKL and surface-area inequalities.
Our recent work \cite{chen2026discrete} also established these stronger inequalities under appropriate assumptions.
\paragraph{Proposed Approach.}
Markov semigroup methods provide a framework for proving influence
and isoperimetric inequalities \cite{CorderoErausquinLedoux12,ODonnellWimmer13,IvanisviliZhang26}.
Our work \cite{chen2026discrete} utilizes the similar idea together with gradient estimates.
We aim to apply the same framework to address the proposed problem.

A \emph{Markov semigroup} $(P_t)_{t\geq 0}$ describes the time evolution of a Markov process and is generated by an operator $L$ ($P_t = e^{tL}$): for any $f: \Omega\mapsto \R$ and $x\in\Omega$, $Lf(x) = \sum_{i=1}^n q_i(x)(f(x^i)-f(x))$.
Here $\{q_i\}_{i\in[n]}$ is called the \emph{transition rates} of the Markov semigroup, which is always assumed to be \emph{reversible} with respect to $\mu$: $\mu(x)q_i(x)=\mu(x^i)q_i(x^i)$.
Given the Markov semigroup $(P_t)_{t\geq 0}$, the corresponding strong gradient estimate holds with a constant $\rho$ if for any function $f:\Omega\mapsto\R$ and $t\geq 0$,
$$
    \norm{\nabla^\star P_t f}_2 \leq e^{-\rho t}P_t\tp{\norm{\nabla^\star f}_2},
$$ 
where the discrete gradient of $f$ is defined as $\nabla^\star f(x) = [q_i(x)(f(x^i)-f(x))]_{i\in[n]}$.
Through the potential function and derivative techniques, the gradient estimate can be equivalently written as a functional inequality, which is known as the reinforced Bakry--\'Emery condition \cite{BakryEmery85}.
This can be formulated as a local matrix eigenvalue problem \cite{CKLP22}.
In \cite{chen2026discrete}, the implication from the gradient estimates to isoperimetric inequalities was established. 
By analyzing commutators between the generator and discrete gradient, we derived sufficient local matrix criteria for gradient estimates and verified them under suitable weak-dependence assumptions on $\mu$.

In our previous work, we chose transition rates $q_i(x) = \frac 12 \sqrt{\mu(x^i)/\mu(x)}$ satisfying the square identity $q_i(x)q_j(x^i) = q_j(x)q_i(x^j)$ and used the commutativity of coordinate flips $(x^i)^j = (x^j)^i$ to establish the commutator identity $(L\nabla^\star - \nabla^\star L)f= H\cdot \nabla^\star f$, where $H: \Omega\mapsto\R^{n\times n}$.
However, for distributions supported on matroid independent sets, coordinate flips may leave the support and the commutativity of elements may not hold, making the square identity fail.
Establishing an appropriate commutator identity is therefore a central challenge.
One potential approach is to redefine gradients and  apply a model-based transition rates.
The directed-derivative identities for birth--death dynamics in
\cite{gobel2026simple} and the local Bochner framework in
\cite{chen2026spectral} suggest a possible route to derive a new commutator identity and establish the gradient estimate.

\bibliographystyle{alpha}
\bibliography{ref.bib}
\end{document}